
We evaluate on TruthfulQA \citep{Lin2021TruthfulQA}, chosen for its adversarial design that challenges UQ methods.

\subsubsection{Dataset}

TruthfulQA (generation task): 817 questions testing factual accuracy on common misconceptions. Human-annotated ground truth. Adversarial design targets known failure modes of language models. Example: ''What happens if you crack your knuckles?'' Data splits: calibration split 326 questions (40\%) for temperature scaling T optimization and conformal threshold calibration; test split 491 questions (60\%) for AUROC evaluation, never seen during calibration. Seed: 42 for reproducible splits.

\subsubsection{Baselines}

\textbf{Temperature Scaling \citep{Guo2017Calibration}} --- Post-hoc logit rescaling. Included as zero-cost baseline, tests whether aleatoric calibration suffices. Cost: 1.$0\times$.

\textbf{Conformal Prediction \citep{Su2024Selective}} --- Distribution-free coverage guarantees. Alternative zero-cost approach with theoretical coverage guarantees. Cost: 1.$0\times$.

\textbf{MC Dropout k=1} --- Deterministic baseline (dropout disabled). Tests whether stochasticity is necessary. Cost: 1.$0\times$.

\textbf{MC Dropout k=3} --- Low-k epistemic uncertainty. Tests minimal k for threshold crossing. Cost: 3.$0\times$.

\textbf{MC Dropout k=5} --- Mid-k Bayesian approximation. Balances cost versus quality. Cost: 5.$0\times$.

\textbf{MC Dropout k=10} --- High-k Bayesian approximation. Tests whether k>5 justified. Cost: 10.$0\times$.

\subsubsection{Implementation Details}

Framework: HuggingFace Transformers (PyTorch backend). Model: Llama-3.1-8B-Instruct. Precision: FP16. Device: auto-mapped to $5\times$ NVIDIA H100 NVL GPUs. Hyperparameters: temperature scaling T optimized via LBFGS on calibration split; conformal prediction $\alpha$ = 0.1 (90th percentile threshold); MC dropout p = 0.1; random seed 42. Compute resources: $5\times$ NVIDIA H100 NVL (80GB each). Total experiment time: ~45 minutes (817 question inference across 6 methods).

\subsubsection{Evaluation Metrics}

\textbf{AUROC:} measures ability to discriminate correct versus incorrect predictions using uncertainty score. Range [0.5, 1.0]. Threshold $\geq$ 0.70 for viable selective prediction. Directly measures selective prediction quality.

\textbf{Inference Cost (FLOPs-normalized):} total FLOPs per query, normalized to single forward pass = 1.$0\times$. MC dropout k=N measured as N.$0\times$. Hardware-agnostic cost metric.

\textbf{Spearman Correlation ($\rho$):} secondary metric measuring monotonic relationship between uncertainty score and incorrectness. Threshold $\rho$ > 0.2 (validation check). Confirms uncertainty scores correlate with prediction correctness.
